\section{Changing the prescribed singularities}

Keep the semipositive big background $\theta$ fixed, and use $B_0,H_0$ for its zero-reference functionals. A common smaller-background condition controls the change of normalization uniformly over all continuous logarithmic densities.

\begin{proposition}\label{prop:reference-normalization-comparison}
Let $\varphi,\psi$ satisfy \eqref{eq:normalization-reference}. If $\varphi\preceq\psi$, then $B_\varphi\leq B_\psi$ and $H_\varphi\leq H_\psi$. Moreover, if $P_{\varepsilon\theta}(\varphi)$ is genuine for $0<\varepsilon<1$, then
\[
 B_0(G)+n\log(1-\varepsilon)
       \leq B_\varphi(G)\leq B_0(G)
       \qquad(G\in C(X)),
\]
and
\[
 H_0(\nu)+n\log(1-\varepsilon)
       \leq H_\varphi(\nu)\leq H_0(\nu)
       \qquad(\nu\in\mathcal P(X)).
\]
The dual inequalities include infinite values.
\end{proposition}
\begin{proof}
Fix $G$, and let $u_\varphi,u_\psi$ solve the two ordinary-type equations, with scalars $c_\varphi,c_\psi$. Then $u_\varphi\preceq u_\psi$. If $c_\varphi>c_\psi$, singular strict domination applied to $u_\psi$ and $u_\varphi+t$ would give $u_\psi\geq u_\varphi+t$ for every real $t$, a contradiction at a common finite point. This proves the monotonicity of $B$, and its defining supremum gives the monotonicity of $H$. Since $\varphi\preceq0$, it also proves both upper bounds.

For the lower bound, let $v$ be the bounded solution for $G$, with scalar $c_0$, and put
\[
 h=P_{\varepsilon\theta}(\varphi)+(1-\varepsilon)v.
\]
The two currents in this sum have backgrounds $\varepsilon\theta$ and $(1-\varepsilon)\theta$, respectively. Positivity and homogeneity give
\[
 \mathrm{MA}^{\mathrm{ext}}_\theta(h)
       \geq(1-\varepsilon)^n\mathrm{MA}_\theta(v)
       =(1-\varepsilon)^nc_0\mathrm{e}^G\,\mathrm{d}V.
\]
As in the lower-density lemma, the inequality is justified first on finite levels and then off the pole set of $P_{\varepsilon\theta}(\varphi)$, which has zero smooth volume. Also $h\preceq u_\varphi$. If $c_\varphi<(1-\varepsilon)^nc_0$, strict domination would force $u_\varphi\geq h+t$ for every $t$. Thus $c_\varphi\geq(1-\varepsilon)^nc_0$. Taking logarithms proves the scalar bound; taking the defining suprema proves the dual bound.
\end{proof}

\begin{corollary}\label{cor:small-reference-adjoint-limits}
Let the model references $\varphi_j$ satisfy $P_{\varepsilon_j\theta}(\varphi_j)\not\equiv-\infty$, where $0<\varepsilon_j<1$ and $\varepsilon_j\to0$. Then $B_{\varphi_j}\to B_0$ uniformly on $C(X)$. If $G_j\to G$ uniformly and $\nu_j\in\mathcal A_{\varphi_j}(G_j)$, every weak cluster point belongs to $\mathcal A_0(G)$. The densities of $\nu_j$ are uniformly integrable, and weak convergence of a subsequence also holds weakly in $L^1(\mathrm{d}V)$. If $\mathcal A_0(G)$ is a singleton, the full sequence converges to it in these two senses.
\end{corollary}
\begin{proof}
Put $\delta_j=-n\log(1-\varepsilon_j)\to0$. The proposition gives $|B_{\varphi_j}-B_0|\leq\delta_j$ on all of $C(X)$. Its dual inequality and the supporting identity give
\[
 0\leq H_0(\nu_j)-\int_XG_j\,\mathrm{d}\nu_j-B_0(G_j)
       \leq\delta_j.
\]
Choose $M$ with $\|G_j\|_\infty\leq M$. This implies $H_0(\nu_j)\leq B_0(0)+2M+\delta_j$, so all $\nu_j$ lie in one compact dual sublevel. If a subsequence converges weakly to $\nu$, lower semicontinuity, uniform convergence of $G_j$ and continuity of $B_0$ give
\[
 H_0(\nu)-\int_XG\,\mathrm{d}\nu-B_0(G)\leq0.
\]
The dual minimum formula gives the opposite inequality and identifies equality with $\nu\in\mathcal A_0(G)$. The common-sublevel theorem \cref{thm:normalization-dual-compactness} also gives uniform integrability and weak $L^1$ convergence. Compactness and uniqueness of the limiting supporting probability give the last assertion.
\end{proof}
